\section{A finite transverse sweep for native actual curl primitives}
\label{sec:r11-finite-native-sweep}

\paragraph{Status, authorship, and the remaining geometric boundary.}
\textbf{PROVED THIS ROUND}, submitted for exact-source independent review.
The construction contributors are \texttt{/root} and
\texttt{/root/r11\_cube}.  This section constructs actual forms on a
specified whole native unit-valuation chart, in every positive carrier
codimension.  It does not assert extension through zero entries of the
original tuple, independence of the chosen unit presentation, or a
global projective primitive.  In particular it does not yet supply the
original-support Green bridge for the general higher-rank cube.

The external definition input is checked against A.~Mihatsch,
\emph{$\delta$-Forms on Lubin--Tate spaces},
\texttt{arXiv:2112.10018v4}, dated 7 September 2024:
Definitions 2.7--2.13, Proposition 2.32 and its proof, and the native
pullback interface of Proposition 2.38 and Lemma 4.4.
The original text was read at
\url{https://arxiv.org/html/2112.10018v4}.
Only the following consequence is used for existence: a finite
polyhedral current on a Euclidean open whose two full derivatives are
polyhedral is a native actual $\delta$-form.  This is not a criterion
for arbitrary currents on arbitrary Berkovich spaces.

\subsection{Fixed finite objects and a transverse rational vector}

Write $V=\mathbb R^d$ and fix a nonzero rational linear coordinate
$\tau$ and the open half-space
\[
                         \Omega=\{\tau>0\}.
 \tag{SWP1}
\]
A finite central carrier presentation on $\Omega$ is the restriction
of finitely many closed polyhedral cones with vertex zero, together
with their weighted integration forms.  Intersections with $\Omega$
are understood relatively; no boundary at $\tau=0$ belongs to our
domain.  Supply an actual form
\[
 F=\sum_{\sigma}\alpha_\sigma\wedge[\sigma,\mu_\sigma]
       \in B^{p+c,q+c}(\Omega),\qquad c\ge1,
       \qquad d'F=d''F=0,
 \tag{SWP2}
\]
where every nonempty carrier has codimension $c$, and
$\alpha_\sigma$ is a constant tangential $(p,q)$-form.  Its value
means a form on $N_\sigma$, not a choice of ambient extension.
The weights $\mu_\sigma\in\det N_\sigma/\{\pm1\}$ are those
of Mihatsch's Definition 2.8; signs of summands may be included in
$\alpha_\sigma$.  Presentations are equivalent by common refinement
and equality of the tangential coefficients with their weights.

There exists a nonzero rational vector
\[
       v\in\ker\tau,\qquad v\notin N_\sigma
                              \quad\hbox{for every carrier.}
 \tag{SWP3}
\]
Indeed, $N_\sigma\cap\ker\tau$ is a proper subspace of
$\ker\tau$.  Otherwise $c=1$, $N_\sigma=\ker\tau$,
and the central carrier would have empty intersection with $\Omega$.
A finite union of proper subspaces cannot contain all rational
vectors in $\ker\tau$.  If $d=1$, the assumptions in (SWP2)
force $F=0$, which has the zero primitive; below assume $d\ge2$.
Fix one vector (SWP3) for the entire given presentation.  All further
subdivisions retain transversality.  No new irrational direction is
needed when the original cones have rational directions.

For each carrier define
\[
 \begin{split}
 \widehat\sigma&=\sigma+\mathbb R_{\ge0}v\subset\Omega,\\
 x&=y+t_\sigma(x)v,
       \quad y\in\sigma,\quad t_\sigma(x)\ge0,\\
 \widehat\mu_\sigma&=\mu_\sigma\wedge v
                     \in\det(N_\sigma+\mathbb Rv)/\{\pm1\}.
 \end{split}
 \tag{SWP4}
\]
The decomposition in (SWP4) is unique.  The function $t_\sigma$
is the restriction of a homogeneous linear function on
$N_\sigma+\mathbb Rv$.  Let $\widehat\alpha_\sigma$ be the
constant $(p,q)$-form on that space obtained from
$\alpha_\sigma$ by projection parallel to $v$.  Thus both of
its contractions by $v$ vanish.  The swept carriers are polyhedral
cones of codimension $c-1$; take a finite common refinement of them
when they overlap.

\begin{theorem}[Finite actual transverse sweep]
The explicit finite polyhedral form
\[
 \boxed{\quad
 \mathcal S_v F=\sum_\sigma
       t_\sigma\widehat\alpha_\sigma
                         \wedge[\widehat\sigma,\widehat\mu_\sigma]
       \in B^{p+c-1,q+c-1}(\Omega)
 \quad}
 \tag{SWP5}
\]
is actual and satisfies
\[
                         D\mathcal S_vF=F,
                         \qquad D=d'd''.
 \tag{SWP6}
\]
For this fixed $v$, it depends only on the original native actual
form, not on its chosen transverse subdivision.  It uses finite
linear coefficient functions, not an infinite-series completion.
\end{theorem}

\begin{proof}
We first check the finite construction as an equality of native
distribution-valued superforms; this will also verify its boundary
coefficients.  For a constant vector $v$, denote its two ordinary
graded contractions on distribution-valued forms by $i'_v,i''_v$.
These individual contractions are not claimed to preserve actual
$\delta$-forms.  Write $T_t(x)=x-tv$, which preserves $\Omega$.
On a compact subset $K\subset\Omega$, the set of parameters
\[
       \{t\ge0:(\sigma+tv)\cap K\ne\varnothing}
 \tag{SWP7}
\]
is bounded, uniformly for the finite list of carriers.  To see this,
choose a linear function $\ell_\sigma$ vanishing on $N_\sigma$
and satisfying $\ell_\sigma(v)=1$.  If $x=y+tv$ belongs to
$K$, then $t=\ell_\sigma(x)$, which is bounded on $K$.
This argument also bounds all terms arising by differentiation or
contraction; these operations do not enlarge the carrier support.
Consequently all the following parameter integrals have compact
parameter support on every compact test support.

A direct calculation in linear coordinates adapted to
$N_\sigma\oplus\mathbb Rv$ gives
\[
 \int_0^\infty t\,T_t^*
           (i''_vi'_v(\alpha_\sigma[\sigma,\mu_\sigma]))\,dt
       =t_\sigma\widehat\alpha_\sigma
                             [\widehat\sigma,\widehat\mu_\sigma].
 \tag{SWP8}
\]
Here is also the normalization and sign check.  Arrange that the
first normal coordinate of $N_\sigma$ has value one on $v$ and
all other normal coordinates vanish on $v$.  The two contractions,
in the order $i''_vi'_v$, remove that prime--double-prime normal
pair.  Moving them past the tangential coefficient produces the
same total-degree sign twice, hence sign $+1$.
Integration in $t$ replaces this normal pair by the swept tangent
direction with weight $\mu_\sigma\wedge v$.
For $F=\delta_0d'x\wedge d''x$ in one dimension and
$v=a\partial_x$, the result is $|x|$ on the ray $x/a\ge0$,
zero on the opposite ray.  Its second derivative is exactly $F$;
the factor $|a|$ in the new weight is essential.
The affine change-of-coordinate rule for weights proves (SWP8)
without any orthonormality or lattice-primitivity assumption.

In the same adapted coordinates the single-contraction formulas are
\[
 \begin{split}
 -\int_0^\infty T_t^*(i''_vF)\,dt
   &=\sum_\sigma d't_\sigma\wedge\widehat\alpha_\sigma
                         [\widehat\sigma,\widehat\mu_\sigma],\\
 \int_0^\infty T_t^*(i'_vF)\,dt
   &=\sum_\sigma d''t_\sigma\wedge\widehat\alpha_\sigma
                         [\widehat\sigma,\widehat\mu_\sigma].
 \end{split}
 \tag{SWP9}
\]
All three right sides (SWP8)--(SWP9) are finite polyhedral forms.
They are not defined only as formal integrals of currents.

The constant-vector Cartan identities on distributions are
\[
 \begin{gathered}
 d'i'_v+i'_vd'=L_v,\qquad d''i''_v+i''_vd''=L_v,\\
 d'i''_v+i''_vd'=0,\qquad d''i'_v+i'_vd''=0,\qquad
                         \partial_tT_t^*=-T_t^*L_v.
 \end{gathered}
 \tag{SWP10}
\]
They follow on coordinate distribution coefficients by the graded
Leibniz rule.  Using $d'F=d''F=0$, (SWP8), and integration by
parts with the bounded support (SWP7), one obtains
\[
 \begin{split}
 d'\mathcal S_vF&=-\int_0^\infty T_t^*(i''_vF)\,dt,\\
 d''\mathcal S_vF&=\int_0^\infty T_t^*(i'_vF)\,dt,\\
 D\mathcal S_vF&=\int_0^\infty T_t^*(L_vF)\,dt=F.
 \end{split}
 \tag{SWP11}
\]
In the first two equations the boundary at $t=0$ is zero because
of the factor $t$.  In the last equation the boundary at $t=0$
is precisely $F$.  The identities use the full closedness of $F$;
its side-face boundary terms cannot be dropped carrier by carrier.

By (SWP9) and (SWP11), the finite form in (SWP5) and both of its
full derivatives are polyhedral.  The native Euclidean case of
Mihatsch's Proposition 2.32 proves that (SWP5) is an actual
$\delta$-form.  Its displayed finite weighted-carrier presentation
is the representative.  Equation (SWP6) is therefore an identity
in native $B$, with every wall coefficient included.

Finally, the left side of (SWP8) is linear in the entire native
current $F$ and is well-defined on every compact support by (SWP7).
Thus any equality of native actual presentations is respected.
Equivalently, subdividing $\sigma$ subdivides its injectively swept
polyhedron, with the same coefficients and weights.  This proves
the asserted refinement compatibility.  No assertion about arbitrary
analytic refinements has been used to create the primitive.
\end{proof}

\subsection{The finite native space needed after the sweep}

The output (SWP5) has an explicit finite polyhedral presentation.
We do not assume, without proof, that it belongs to the smaller
algebra of differential polynomials in the original PL generators
from CE1.  Instead the following elementary extension supplies the
precise operator domain needed for the two-colour solver.

Let $\mathscr P_{\rm cone}(\Omega)$ consist of native actual
$\delta$-forms with a finite central polyhedral presentation and
polynomial tangential coefficients.  We allow a finite subdivision
of its carriers.  Equivalence is equality of those actual forms on
a common refinement.  All derivative operations are the full native
derivatives.  Polynomial coefficients, their tangential derivatives,
and all facet traces remain polynomial on a finite central
refinement, so $d',d''$ preserve this space.

Let $E=\sum_jx_j\partial_{x_j}$.  Define $I'$ and $I''$
on a presented form by contraction of its tangential coefficients
by $E|_\sigma$.  Define $N$ on tangential coefficients by sending
every prime covector to its double-prime counterpart and killing
double-prime covectors.  The radial vector is tangent to every
central carrier and all of its faces.  These operations therefore
respect common refinement and actual equality of coefficients.
They give finite polyhedral forms before any actualness assertion.

On a homogeneous polynomial coefficient of degree $m$ and
tangential bidegree $(p,q)$, set
\[
                         E'=m+p,\qquad E''=m+q.
 \tag{SWP12}
\]
These numbers do not count the codimension of the carrier.  As
distribution-valued forms the exact identities are
\[
 \begin{gathered}
 d'I'+I'd'=E',\qquad d''I''+I''d''=E'',\\
 d''I'+I'd''=N,\qquad d'I''+I''d'=N^+,\\
 [N,d']=d'',\qquad [N,d'']=0.
 \end{gathered}
 \tag{SWP13}
\]
For completeness these identities may be checked in linear normal
coordinates for each carrier.  Radial contraction kills every normal
pair on the carrier, and $N$ kills that pair because its image has a
repeated double-prime normal factor.  The normal distribution is
homogeneous of degree minus its codimension, cancelling the normal
covector contribution to either Euler weight.  What remains is
exactly the tangential coefficient calculation (SWP12).  The same
calculation on all faces gives the distributional identities,
including the facet derivatives.  Alternatively (SWP13) follows
directly from the coordinate Cartan identities on distributions
and the just-described homogeneous normal factors.

For $T\in\mathscr P_{\rm cone}(\Omega)$, both
\[
 d'I'T=E'T-I'd'T,\qquad d''I'T=NT-I'd''T
 \tag{SWP14}
\]
are finite polyhedral forms; the analogous assertions hold for
$I''T$.  Proposition 2.32 again proves $I'T,I''T$ actual.
The commutators in (SWP13) similarly prove $NT$ actual.  Thus the
operators in (SWP13) act on this explicitly specified finite native
space, with no use of an abstract quotient inverse.  They agree
with CE2 on the original finite generator algebra.

For a fixed carrier complex, a fixed vector $v$, and a fixed bound
on polynomial degree, all constructions above are linear maps
between finite-dimensional coefficient spaces and hence continuous.
The target refinement is fixed for the whole input coefficient
space, not enlarged coefficient by coefficient.  This statement
does not identify any whole Berkovich coefficient-final topology,
give a universal continuous selector, or solve OPEN-KERNEL.

\subsection{An explicit native third witness in every positive codimension}

Supply a triple $u=(A,B,M)$ on $\Omega$ with total bidegrees
\[
 (c,c+2),\qquad(c+2,c),\qquad(c+1,c+1),\qquad c\ge1,
 \tag{SWP15}
\]
lying in $\mathscr P_{\rm cone}(\Omega)$, with Euler weights
$(1,3),(3,1),(2,2)$ and
\[
                 d'A+d''M=0,\qquad d''B+d'M=0.
 \tag{SWP16}
\]
Assume its curls $C=d''A$ and $C^+=d'B$ have finite
constant-coefficient central presentations of pure carrier
codimension $c$, tangential bidegrees $(0,3)$ and $(3,0)$.
This is a finite, entire-presentation condition, not an assertion
that arbitrary actual curls have this shape.
Both curls are independently closed by (SWP16) and
anticommutation.  Choose a single rational $v$ transverse to the
union of both finite carrier lists and put
\[
 F=\mathcal S_vC,\qquad G=\mathcal S_vC^+,\qquad
 A^\circ=A+d'F,\quad B^\circ=B-d''G,\quad M^\circ=M.
 \tag{SWP17}
\]
Then $DF=C$, $DG=C^+$, the corrected triple is
$\partial_2$-closed, and
\[
                        d''A^\circ=d'B^\circ=0.
 \tag{SWP18}
\]
The swept coefficients have degree one, so their Euler weights are
$(1,4)$ and $(4,1)$ respectively; the corrected triple retains
the original three weights.

Define the following explicit finite forms in the enlarged native
space just verified:
\[
 \begin{gathered}
 P_0=\tfrac13 I''A^\circ,\qquad Q_0=\tfrac13 I'B^\circ,\\
 M_0=M-d'P_0-d''Q_0,\qquad P_1=\tfrac12 I'M_0,\qquad
 U_1=-\tfrac12 I'NM_0.
 \end{gathered}
 \tag{SWP19}
\]
The literal identity is
\[
 \boxed{\quad
 u=\partial_3(U_1-F, G, P_0+P_1, Q_0).
 \quad}
 \tag{SWP20}
\]
Indeed (SWP13) and (SWP18) give $d''P_0=A^\circ$ and
$d'Q_0=B^\circ$.  Then $d'M_0=d''M_0=0$,
$d'P_1=M_0$, and $d''P_1=NM_0/2$.
The closed form $NM_0$ has prime Euler weight one, so
$d'U_1=-NM_0/2$.  This proves the three components of
(SWP20), with the remaining correction signs from (SWP17).
This is the CGS algebraic argument, now with the previously missing
native operator-domain verification after a sweep; it does not
silently identify two different finite algebras.

\subsection{Analytic substitution and the exact unresolved extension}

Let $\mathcal V$ be a pure-dimensional original-$K$ analytic space
without boundary, supplied with an entire finite tuple of analytic
units whose valuation map takes values in $\Omega$.  The native
forms (SWP5) and (SWP19), as well as their identities, pull back to
actual Mihatsch forms on $\mathcal V$ by the original native
presentation interface.  Every refinement and all mixed coefficients
are retained: the identity is in the full Euclidean $B(\Omega)$,
not merely on the image of a selected skeleton or on probe products.
There is no extension of arbitrary-current pullback in this step.
One may specialize the auxiliary unit to an original element
$\pi\in K$ with $v(\pi)>0$ after this actual construction.

For the native universal raw capped cube, the finite curvature
products and first valuation differentials give exactly the pure
constant-coefficient curl presentations required above.  Thus (SWP20)
is available on its full unit-log chart once its full closedness and
weight identities have been read from the exact CS/MT formulas.
This statement has no conclusion at a proper-entry-zero point of
the original tuple.  A swept coefficient can depend on an individual
unit logarithm that becomes infinite there.  The finite norm/cap
preservation lemma NCP does not apply to such a coefficient merely
because its curvature equation holds on the unit locus.

In particular, extending these native primitives across the original
proper-entry zero sets, making them compatible across projective
charts, and realizing the resulting higher-rank witness on the
original exact $U$ remain \textbf{OPEN}.  No original cycle is
changed, but no general original-cycle Green form is claimed by this
native construction alone.  The known strict proper-fiber
obstructions are consistent with a strict primitive on an affine
unit-valuation chart.  The whole analytic OPEN-KERNEL and its
inherited topology likewise remain open.

\paragraph{Exact internal sources.}
CE: \path{68a14ad23c956c211aaae0c188aebb672332181b0a13a269e21cda595a0b6f1d}.
CGS: \path{dc64c1fb85dbd8dfba2105a230e11dc563189df38a4d681e587b3d5784fc6d21}.
CS: \path{d0c06d21c2289d292fb735d6966d456b9e05f1e4fab10f263a5007169d79cacf}.
MT: \path{d1365f528071a8db5e20f55628c88b84a06e73965dd808e8b27375c776fc73d4}.
NCP: \path{144df3e33ec6beeefa7907b1292edc31855d289727da777082c5015d33676743}.
This section does not replace or alter any of these frozen sources.
